\label{chap:p3-83-25-alpha}

\section[Internal derivation of \texorpdfstring{\(\alpha_{\rm HMT}\)}{alpha HMT}]
{Internal derivation of \texorpdfstring{\(\alpha_{\rm HMT}\)}{alpha HMT}:
common terminal state, dodecaphase seal, and carry cocycle}
\label{p3-83-c25-s1:chap:parte-iii-alpha-hmt}

The preceding generator reaches a terminal state enriched with two
coordinated publications:
\[
\begin{aligned}
w_6\longrightarrow w_{30}\longrightarrow R_{36}\longrightarrow G_9
\longrightarrow x_{\rm term}
&\xrightarrow{\ \pi_{\rm term}\ }B_{\rm term},\\
x_{\rm term}
&\xrightarrow{\ R_{\rm sgn}\ }U_{12}^{\rm sgn}
\xleftrightarrow{\ \mathcal H_{12}\ }K
\longrightarrow A\longrightarrow\alpha_{12}.
\end{aligned}
\]
The two terminal arrows publish distinct coordinates of the same antecedent.
The transform \(\mathcal H_{12}\) is reversible; the transverse system
\((D_3K,D_4K,Q(K))\) attains rank twelve; oriented Euclidean division preserves
the carry cocycle and the boundary conditions. The chain uses the
signed state, the seal, and the declared primitives, and obtains
\(\alpha_{\rm HMT}\) before opening metrological recognition.

\subsection{Dodecaphasic closure and compatible constructions}
Dodecaphase closure provides the governing statement and proof. The
E21/22 constructions and \(\mathsf T_{\alpha,\infty}\) retain their own domains and
functions within the principal genealogy.


\section{Signed register, seal reconstruction, and integral carry}
The signed register and the carry cocycle unfold the common terminal state,
the Hadamard charts and transverse differences, the Smith normal form,
the affine identity, the coinductive prolongation, and the boundaries
\(662/663\).


\section{Reversibility, independence, and non-circularity of the dodecaphasic closure}
The state \(108\), the prospective constructor of \(K\), the bidirectional
Hadamard transform, integral carry, and infinite prolongation are
articulated through explicit controls of reversibility, independence, and
noncircularity. Each formulation preserves its domain and its
refutation criterion.


% Unidad U016. Registro firmado, lectura terminal y reconstruccion de K.
% Redaccion autonoma desde la autoridad material 1063 y su sucesor V2.

\section{Factorization of the signed dodecaphasic record and integral closure of the seal \(K\)}
\label{p3-83-c25-s4:chap:registro-dodecafasico-hadamard-k}

The construction of the integer seal does not begin with \(\alpha\), its inverse
or a metrological table. Its domain is the enriched terminal state
of the Topological Phase Kernel obtained after the already constructed chain
\begin{equation}
 \boxed{
 \mathrm{APP}\longrightarrow\mathrm{TRIT}\longrightarrow\mathrm{TPK}
 \longrightarrow\Omega_{\rm seed}
 \xrightarrow{T_{\min}}\mathcal U_6^{\rm cat}
 \xrightarrow{\rho_3}W_{243}
 \longrightarrow w_{30}\longrightarrow R_{36}
 \longrightarrow G_9\longrightarrow x_{\rm term}.}
 \label{p3-83-c25-s4:eq:cadena-upstream-registro-dodecafasico}
\end{equation}
The foregoing arrows were defined with their respective domains:
pairs of seeds, signatures, emissions, ternary words, boundaries,
compatible histories, and nonadic holonomy. This unit does not reselect
any of those data. It constructs the integral charts of the same
terminal state and proves their equivalence.

\subsection{Windows of the oriented trace of length \(108\)}

Let
\[
 \Theta_{108}=\mathbb Z/108\mathbb Z
\]
the observable calendar. Its twelve nonadic windows are
\[
 E_m=\{9m+1,\ldots,9m+9\},
 \qquad 0\leq m<12.
\]
The partition is disjoint and exhaustive:
\[
 \{1,\ldots,108\}=\bigsqcup_{m=0}^{11}E_m.
\]
The set of address codes that publish an event is fixed as
\[
 D_{\rm evt}=\{2,4,6,8\}.
\]
The orientation of the twelve windows is
\begin{equation}
 \Sigma_{12}=(+,+,+,-,-,-,+,+,+,-,-,-).
 \label{p3-83-c25-s4:eq:signatura-doce-ventanas}
\end{equation}
If \(d_t\) is the direction encoded at step \(t\), the signed event
of window \(m\) is
\begin{equation}
 e_m=\Sigma_{12}(m)\,
 \#\{t\in E_m:d_t\in D_{\rm evt}\}.
 \label{p3-83-c25-s4:eq:evento-firmado-ventana}
\end{equation}
The formula distinguishes the nonnegative number of events from the sign of the
orientation. It does not replace a cardinal direction with a sign: both
coordinates remain present in the TPK state.

\begin{definition}[Enhanced dodecaphonic registration]
\label{p3-83-c25-s4:def:registro-enriquecido-dodecafasico}
The record transported by the trace is
\[
 \mathcal R_{12}(x)=
 \bigl((E_m,\tau_m,\sigma_m,\kappa_m,\Lambda_m)\bigr)_{m=0}^{11},
\]
where \(\tau_m\) is the TRIT orientation, \(\sigma_m\) the boundary
signature, \(\kappa_m\) the carry, and \(\Lambda_m\) the genealogical register segment
that preserves the route. There is a chain of projections
\[
 \mathcal X_{\rm TPK}^{\rm enr}
 \longrightarrow\mathcal R_{12}
 \longrightarrow\Theta_{108},
\]
but neither of the two arrows is an identification.
\end{definition}

In particular, \(\Theta_{108}\) preserves only the position in the calendar.
The record \(\mathcal R_{12}\) additionally preserves orientation, signature,
carry, and provenance. The complete state further preserves cylinder,
boundary, quotient, terminal, and vertical memory.

\subsection{Integral chart of the dodecaphasic opposition}

For \(0\leq i<6\), the opposite positions determine
\begin{equation}
 p_i=e_i+e_{i+6},
 \qquad
 a_i=e_i-e_{i+6}.
 \label{p3-83-c25-s4:eq:pares-opuestos-registro}
\end{equation}
The central charge and the five relative margins are
\begin{equation}
 s_C=\sum_{i=0}^{5}p_i,
 \qquad
 M_i=p_i-p_5\quad(0\leq i<5).
 \label{p3-83-c25-s4:eq:carga-margenes-registro}
\end{equation}

\begin{definition}[Integral chart \(1+5+6\)]
\label{p3-83-c25-s4:def:carta-integral-156}
It is defined
\begin{equation}
 \mathcal L_{12}(e_0,\ldots,e_{11})
 =\bigl(s_C,M_0,\ldots,M_4,a_0,\ldots,a_5\bigr).
 \label{p3-83-c25-s4:eq:carta-integral-156}
\end{equation}
The coordinate partition is
\[
 1+5+6=12.
\]
\end{definition}

\begin{theorem}[Exact inverse of the integral chart]
\label{p3-83-c25-s4:thm:inversion-carta-integral-156}
A tuple
\[
 \ell=(s_C,M_0,\ldots,M_4,a_0,\ldots,a_5)\in\mathbb Z^{12}
\]
belongs to the image of \(\mathcal L_{12}\) if and only if
\begin{align}
 s_C-\sum_{i=0}^{4}M_i&\equiv0\pmod6,
 \label{p3-83-c25-s4:eq:congruencia-seis-carta}\\
 p_i&\equiv a_i\pmod2
 \quad(0\leq i<6),
 \label{p3-83-c25-s4:eq:congruencia-paridad-carta}
\end{align}
where
\begin{align}
 p_5&=\frac{s_C-\sum_{i=0}^{4}M_i}{6},
 \label{p3-83-c25-s4:eq:inversa-p5-carta}\\
 p_i&=p_5+M_i\quad(0\leq i<5).
 \label{p3-83-c25-s4:eq:inversa-pi-carta}
\end{align}
In that case, the only entire preimage is
\begin{equation}
 e_i=\frac{p_i+a_i}{2},
 \qquad
 e_{i+6}=\frac{p_i-a_i}{2}.
 \label{p3-83-c25-s4:eq:inversa-eventos-carta}
\end{equation}
\end{theorem}

\begin{proof}
Summing \(M_i=p_i-p_5\) for \(0\leq i<5\), we obtain
\[
 \sum_{i=0}^{4}M_i
 =\sum_{i=0}^{4}p_i-5p_5
 =s_C-6p_5,
\]
which forces \eqref{p3-83-c25-s4:eq:inversa-p5-carta} and the congruence modulo six.
Once the six \(p_i\) have been reconstructed, the system
\[
 p_i=e_i+e_{i+6},
 \qquad a_i=e_i-e_{i+6}
\]
has the unique solution \eqref{p3-83-c25-s4:eq:inversa-eventos-carta}. This solution is
integral exactly when \(p_i\) and \(a_i\) have the same parity.
The two families of congruences are therefore necessary and sufficient.
\end{proof}

This theorem precedes the numerical closure of \(\alpha\). It explains how the
oriented trace produces a reversible integral chart, and why preserving
the reduced calendar is not enough.

\subsection{Selection of the common terminal state}

The three compatible sections reach terminal depth with catalogs
of cardinalities
\[
 |\mathcal T_\pi|=196,
 \qquad
 |\mathcal T_e|=197,
 \qquad
 |\mathcal T_\varphi|=196.
\]
The product contains
\begin{equation}
 196\cdot197\cdot196=7\,567\,952
 \label{p3-83-c25-s4:eq:cardinal-producto-terminal}
\end{equation}
ordered triples.

The column margin of the terminal reader is
\[
 q_\partial=(1,2,2,1,2,0)\in\mathbb F_3^6,
\]
and the row orientation of the combined channel is
\[
 \sigma=(1,1,-1)=(1,1,2)\in\mathbb F_3^3.
\]
The first datum acts on columns; the second acts on rows. Neither is deducible from the other.

The local signatures of Gram matrix, norm, weight, and distance reduce
\eqref{p3-83-c25-s4:eq:cardinal-producto-terminal} to \(331\) triples. The margin
\(q_\partial\) leaves two indices:
\[
 (120,197,157),
 \qquad
 (129,197,148).
\]
Their matrices are
\[
 B_0=
 \begin{pmatrix}
 0&2&0&2&2&0\\
 0&0&1&2&2&0\\
 1&0&1&0&1&0
 \end{pmatrix},
 \qquad
 B_1=
 \begin{pmatrix}
 0&2&1&0&2&0\\
 0&0&1&2&2&0\\
 1&0&0&2&1&0
 \end{pmatrix}.
\]
The row sums are
\[
 (0,2,0),
 \qquad
 (2,2,1).
\]
Since
\[
 \langle\sigma\rangle
 =\{(0,0,0),(1,1,2),(2,2,1)\},
\]
only \(B_1\) satisfies the axial condition.

\begin{theorem}[Uniqueness of the Visible Terminal Publication]
\label{p3-83-c25-s4:thm:unicidad-publicacion-terminal-visible}
The joint selection determines
\begin{equation}
 B_{\rm term}=
 \begin{pmatrix}
 021020\\001220\\100210
 \end{pmatrix},
 \label{p3-83-c25-s4:eq:matriz-terminal-visible}
\end{equation}
corresponding to the triple \((129,197,148)\). No decimal evaluation of
\(\alpha\) enters the census or the two selection conditions.
\end{theorem}

\begin{proof}
The census and margin leave exactly \(B_0,B_1\). The charge of \(B_0\)
does not belong to \(\langle\sigma\rangle\), whereas that of \(B_1\) is
\(2\sigma\). The second candidate is therefore the only one that simultaneously satisfies
both typed requirements.
\end{proof}

\subsection{2 publications from the same state and absence of retroactive conversion}

Let
\[
 x_{\rm term}\in
 \mathcal X_{\rm TPK}^{\rm term,enr}
 \subseteq\mathcal X_{\rm TPK}^{[108]}.
\]
The visible reading is
\[
 \pi_{\rm term}(x_{\rm term})=B_{\rm term}.
\]
The signed chart preserves another coordinate of the same state.

\begin{definition}[Terminal signed causal reading]
\label{p3-83-c25-s4:def:subespacio-firmado-terminal}
The domain incorporates neither \(K\) nor \(U\) as input coordinates. The
state \(x_{\rm term}\), produced by the full composition
\(\mathrm{APP}\to\mathrm{TRIT}\to\mathrm{TPK}\), contains the enriched
terminal register, and the extractor of its twelve oriented windows is
\begin{equation}
 R_{12}:\mathcal X_{\rm TPK}^{\rm term,enr}
 \longrightarrow\mathcal R_{12}^{\rm or},
 \qquad
 R_{12}(x_{\rm term})
 =\bigl((E_m,\tau_m,\sigma_m,\kappa_m,\Lambda_m)\bigr)_{m=0}^{11}.
 \label{p3-83-c25-s4:eq:subespacio-firmado-terminal}
\end{equation}
The ledger’s channel reader \(90/120\) emits, without receiving the result
as data,
\[
 E_{108}^{90,120}:\mathcal R_{12}^{\rm or}
 \longrightarrow\mathbb Z^{12}\times\mathbb Z^{12}\times\mathbb Z,
 \qquad
 E_{108}^{90,120}\bigl(R_{12}(x_{\rm term})\bigr)
 =(b_{90},b_{120},Q_{\rm TPK}),
\]
with
\begin{align*}
 b_{90}={}&(-495,-116,-684,108,601,-90,
             -173,-88,313,560,-397,461),\\
 b_{120}={}&(-425,-281,-481,671,-255,30,
              475,-543,680,251,6,-128),\\
 Q_{\rm TPK}={}&6263.
\end{align*}
For \(r=1,2,3\) and \(j=0,1,2,3\), define
\[
\begin{aligned}
 B_r&=\sum_{j=0}^{3}b_{120,r+3j},
 &d_{r,j}&=b_{90,r+3j},\\
 A_1&=\frac{Q_{\rm TPK}+2B_1+B_2}{3},
 &A_2&=A_1-B_1,
 &A_3&=A_2-B_2.
\end{aligned}
\]
The publication \(\Pi_H\) forms the three blocks
\[
 U^{(r)}=
 \bigl(A_r,d_{r,1}-d_{r,3},d_{r,0}+d_{r,2},
 d_{r,0}-d_{r,2}\bigr),
\]
and reinterleaves them in dodecaphasic order, that is,
\[
 \Pi_H(b_{90},b_{120},Q_{\rm TPK})
 :=\operatorname{reint}_3\bigl(U^{(1)},U^{(2)},U^{(3)}\bigr)=U.
\]
Let \(\mathcal H_{12}\) be the
integral transformation to be constructed in the next section. The publication
module is
\begin{equation}
 \mathcal U_{12}^{\rm sgn}
 :=\operatorname{im}
 \bigl(\mathcal H_{12}:\mathbb Z^{12}\to\mathbb Z^{12}\bigr).
 \label{p3-83-c25-s4:eq:modulo-publicacion-firmada}
\end{equation}
The signed reading is, causally, the composition
\begin{equation}
 R_{\rm sgn}:=
 \Pi_H\circ E_{108}^{90,120}\circ R_{12}:
 \mathcal X_{\rm TPK}^{\rm term,enr}
 \longrightarrow\mathcal U_{12}^{\rm sgn},
 \qquad
 R_{\rm sgn}(x_{\rm term})=U.
 \label{p3-83-c25-s4:eq:lector-firmado-terminal}
\end{equation}
Only after this output is
\(K=\mathcal H_{12}^{-1}U\) reconstructed; the identity
\(U=\mathcal H_{12}K\) certifies that reconstruction and does not define the reader’s
domain.
\end{definition}

The correct joint writing is
\begin{equation}
 x_{\rm term}
 \xrightarrow{(\pi_{\rm term},R_{\rm sgn})}
 (B_{\rm term},U_{12}^{\rm sgn}).
 \label{p3-83-c25-s4:eq:doble-publicacion-terminal}
\end{equation}
No arrow
\(B_{\rm term}\to U_{12}^{\rm sgn}\) is interposed: the visible matrix has forgotten
precisely the orientation, opposition, quotient, carry, and genealogical
record coordinates used by the signed chart. The source of both publications is
the enriched state produced by \eqref{p3-83-c25-s4:eq:cadena-upstream-registro-dodecafasico}.

The evaluation of the signed record gives
\begin{equation}
 \boxed{
 U=(2378,1406,2479,-452,998,-551,
 -668,-204,-371,-322,-28,-997).}
 \label{p3-83-c25-s4:eq:vector-u-firmado-autonomo}
\end{equation}

\subsection{Naturality with respect to depth}

For \(N'\geq N\geq108\), the truncations of the signed system are
\begin{equation}
 \tau_{N',N}^{\rm sgn}
 x_{\rm term}^{[N']}
 =\tau_{N',N}x_{\rm term}^{[N']}.
 \label{p3-83-c25-s4:eq:truncamiento-preserva-ku-autonomo}
\end{equation}
Therefore, truncation acts only on the history and the enriched
state; it does not carry either \(K\) or \(U\) as independent inputs.

\begin{theorem}[Square of naturality of the signed reading]
\label{p3-83-c25-s4:thm:naturalidad-lector-firmado-autonomo}
We have
\begin{equation}
 \boxed{
 R_{{\rm sgn},N}\circ\tau_{N',N}^{\rm sgn}
 =\operatorname{id}_{\mathcal U_{12}^{\rm sgn}}
  \circ R_{{\rm sgn},N'}.}
 \label{p3-83-c25-s4:eq:cuadrado-naturalidad-lector-firmado}
\end{equation}
\end{theorem}

\begin{proof}
The naturality of the states literally preserves the first one hundred and eight
windows of the ledger:
\[
 R_{12,N}\circ\tau_{N',N}^{\rm sgn}=R_{12,N'}.
\]
Truncation can modify only the subsequent history. Since
\(E_{108}^{90,120}\) and \(\Pi_H\) depend exclusively on those windows,
both sides return the same \(U\). Commutativity is an equality of
maps, not an analogy between diagrams.
\end{proof}

\subsection{Orbits of step \(3\) and transform of Hadamard}

The cyclic order of the seal is distributed over the three orbits.
\begin{align}
 K^{(1)}&=(K_1,K_4,K_7,K_{10}),
 \label{p3-83-c25-s4:eq:orbita-k-1}\\
 K^{(2)}&=(K_2,K_5,K_8,K_{11}),
 \label{p3-83-c25-s4:eq:orbita-k-2}\\
 K^{(3)}&=(K_3,K_6,K_9,K_{12}).
 \label{p3-83-c25-s4:eq:orbita-k-3}
\end{align}
Let
\begin{equation}
 H_4=
 \begin{pmatrix}
 1&1&1&1\\
 1&1&-1&-1\\
 1&-1&1&-1\\
 1&-1&-1&1
 \end{pmatrix}.
 \label{p3-83-c25-s4:eq:matriz-hadamard-cuatro-autonoma}
\end{equation}
If \(P_{\rm orb}\) reorders the twelve coordinates according to
\eqref{p3-83-c25-s4:eq:orbita-k-1}--\eqref{p3-83-c25-s4:eq:orbita-k-3}, the global transformation is
\begin{equation}
 \mathcal H_{12}
 =P_{\rm orb}^{-1}
 (H_4\oplus H_4\oplus H_4)P_{\rm orb}.
 \label{p3-83-c25-s4:eq:hadamard-global-doce}
\end{equation}

The three blocks of \eqref{p3-83-c25-s4:eq:vector-u-firmado-autonomo}, read by
orbits, are
\begin{align*}
 U^{(1)}&=(2378,-452,-668,-322),\\
 U^{(2)}&=(1406,998,-204,-28),\\
 U^{(3)}&=(2479,-551,-371,-997).
\end{align*}

\begin{lemma}[Algebraic Hadamard quarter-turn]
\label{p3-83-c25-s4:lem:cuarto-inversa-hadamard}
The matrix \(H_4\) satisfies
\begin{equation}
 H_4^2=4I_4,
 \qquad
 H_4^{-1}=\frac14H_4,
 \qquad
 \det H_4=-16.
 \label{p3-83-c25-s4:eq:identidades-hadamard-cuatro}
\end{equation}
\end{lemma}

\begin{proof}
The rows of \(H_4\) are orthogonal, and each has squared norm four.
Since the matrix is symmetric, \(H_4^{\mathsf T}H_4=H_4^2=4I_4\). The
formula for the inverse and the determinant follow immediately.
\end{proof}

\begin{theorem}[Unique full reconstruction of the seal]
\label{p3-83-c25-s4:thm:reconstruccion-entera-k-hadamard}
The inverse
\begin{equation}
 K^{(r)}=\frac14H_4U^{(r)}
 \qquad(r=1,2,3)
 \label{p3-83-c25-s4:eq:inversion-bloques-hadamard}
\end{equation}
yields
\begin{align*}
 K^{(1)}&=(234,729,621,794),\\
 K^{(2)}&=(543,659,058,146),\\
 K^{(3)}&=(140,824,914,601).
\end{align*}
By undoing the orbital order one obtains
\begin{equation}
 \boxed{
 K=(234,543,140,729,659,824,621,058,914,794,146,601).}
 \label{p3-83-c25-s4:eq:sello-k-doce-autonomo}
\end{equation}
It is the unique rational preimage of \(U\), and it belongs to
\(\mathbb Z^{12}\).
\end{theorem}

\begin{proof}
By \eqref{p3-83-c25-s4:eq:identidades-hadamard-cuatro}, the rational preimage is unique. The explicit multiplications are
\begin{align*}
 H_4U^{(1)}&=(936,2916,2484,3176),\\
 H_4U^{(2)}&=(2172,2636,232,584),\\
 H_4U^{(3)}&=(560,3296,3656,2404).
\end{align*}
All coordinates are divisible by four. Dividing and reinterleaving
yields \eqref{p3-83-c25-s4:eq:sello-k-doce-autonomo}.
\end{proof}

The factor \(1/4\) is not chosen to fit an output: it is the inverse
forced by \(H_4^2=4I_4\). Nor is it the square root of \(-1\). The
internal complex structure \(A_W^2=-I\) will appear after constructing the
Paley matrix; these are two facts of different types.

\subsection{Integral structure and Smith normal form}

\begin{theorem}[Integral image of \(H_4\)]
\label{p3-83-c25-s4:thm:smith-hadamard-autonomo}
The Smith normal form is
\begin{equation}
 \operatorname{SNF}(H_4)=\operatorname{diag}(1,2,2,4).
 \label{p3-83-c25-s4:eq:snf-hadamard-autonoma}
\end{equation}
Consequently,
\begin{equation}
 \operatorname{coker}\mathcal H_{12}
 \cong(\mathbb Z/2\mathbb Z)^6
 \oplus(\mathbb Z/4\mathbb Z)^3,
 \label{p3-83-c25-s4:eq:cociente-hadamard-doce}
\end{equation}
and the image of \(\mathcal H_{12}\) has index
\[
 16^3=4096
\]
in \(\mathbb Z^{12}\). For \(u\in\mathbb Z^4\),
\begin{equation}
 u\in H_4\mathbb Z^4
 \quad\Longleftrightarrow\quad
 H_4u\equiv0\pmod4.
 \label{p3-83-c25-s4:eq:criterio-imagen-hadamard}
\end{equation}
\end{theorem}

\begin{proof}
The greatest common divisor of the entries is one. The minors of order two are
even, and one has absolute value two; those of order three are multiples of
four, and one has absolute value four; the determinant has modulus
sixteen. The determinantal divisors are \(1,2,4,16\), yielding
the invariant factors \(1,2,2,4\). The triple direct sum repeats the
factors and multiplies the indices. Finally,
\(H_4^{-1}=H_4/4\), so the preimage is integral exactly under the
congruence \eqref{p3-83-c25-s4:eq:criterio-imagen-hadamard}.
\end{proof}

\subsection{Certification and reconstruction by differences of steps \(3\) and \(4\)}

Once the ledger has emitted \((b_{90},b_{120},Q_{\rm TPK})\),
\(\Pi_H\) has produced \(U\) and Hadamard has reconstructed \(K\), the
following coordinates act downstream as an independent certificate of
those prior outputs. With cyclic indices modulo twelve, define
\begin{equation}
 (D_3K)_m=K_m-K_{m+3},
 \qquad
 (D_4K)_m=K_m-K_{m+4},
 \qquad
 Q(K)=\sum_{m=1}^{12}K_m.
 \label{p3-83-c25-s4:eq:extractor-diferencias-k}
\end{equation}
For \eqref{p3-83-c25-s4:eq:sello-k-doce-autonomo}, one verifies
\begin{align*}
 D_3K=b_{90}={}&(-495,-116,-684,108,601,-90,\\
         &\qquad-173,-88,313,560,-397,461),\\
 D_4K=b_{120}={}&(-425,-281,-481,671,-255,30,\\
         &\qquad475,-543,680,251,6,-128),\\
 Q(K)=Q_{\rm TPK}={}&6263.
\end{align*}

\begin{theorem}[Completeness of the transversal extractor]
\label{p3-83-c25-s4:thm:completitud-extractor-transversal-k}
The operator
\[
 \mathcal E_{\rm dod}=(D_3,D_4,Q):
 \mathbb Z^{12}\longrightarrow\mathbb Z^{25}
\]
has rank twelve and Smith normal form
\begin{equation}
 \operatorname{SNF}(\mathcal E_{\rm dod})
 =\operatorname{diag}(\underbrace{1,\ldots,1}_{11},12).
 \label{p3-83-c25-s4:eq:snf-extractor-transversal-k}
\end{equation}
In particular, \((D_3K,D_4K,Q(K))\) determines a unique \(K\). This
uniqueness certifies and reconstructs the data already emitted by the ledger; it neither
generates nor selects them.
\end{theorem}

\begin{proof}
If \(D_3v=D_4v=0\), then \(v\) is constant along steps
three and four. Since \(\gcd(3,4)=1\), those steps generate all of
\(\mathbb Z/12\mathbb Z\), so that
\[
 \ker D_3\cap\ker D_4=\langle\mathbf1\rangle.
\]
The total charge does not vanish on that line, since \(Q(\mathbf1)=12\). The rank
is twelve. The rows of \((D_3,D_4)\) are oriented incidences of a connected
Cayley graph; a spanning tree supplies eleven unit
invariant factors. Every maximal minor must contain the row \(Q\), and the
unimodular operations reduce its final contribution to twelve. There is a maximal minor of
modulus twelve; therefore, the last factor is exactly twelve.
\end{proof}

The certifying reconstruction of \(U\) from these coordinates reproduces the
prior publication of \(\Pi_H\) and does not require multiplying again by
\(H_4\). For \(r=1,2,3\), let
\begin{equation}
 B_r=\sum_{j=0}^{3}(D_4K)_{r+3j}
 =\sum_{j=0}^{3}b_{120,r+3j}.
 \label{p3-83-c25-s4:eq:br-diferencias-k}
\end{equation}
Then
\[
 (B_1,B_2,B_3)=(972,-1073,101).
\]
Orbital sums are
\begin{align}
 A_1&=\frac{Q+2B_1+B_2}{3}=2378,
 \label{p3-83-c25-s4:eq:a1-reconstruccion-u}\\
 A_2&=A_1-B_1=1406,
 \label{p3-83-c25-s4:eq:a2-reconstruccion-u}\\
 A_3&=A_2-B_2=2479.
 \label{p3-83-c25-s4:eq:a3-reconstruccion-u}
\end{align}
If \(d_{r,j}=(D_3K)_{r+3j}=b_{90,r+3j}\), the signed block of orbit
\(r\) is
\begin{equation}
 \bigl(A_r,
 d_{r,1}-d_{r,3},
 d_{r,0}+d_{r,2},
 d_{r,0}-d_{r,2}\bigr),
 \label{p3-83-c25-s4:eq:bloque-firmado-desde-diferencias}
\end{equation}
which reproduces the three blocks \(U^{(r)}\) already published by \(\Pi_H\).

\subsection{Threshold Tetrad of the Seal}

The local decimal completion satisfies
\[
 1000=729+271.
\]
The support of the seal coordinates that reach the threshold \(729\) is
\begin{equation}
 \boxed{
 B_K:=\{m:K_m\geq729\}=\{4,6,9,10\}.}
 \label{p3-83-c25-s4:eq:tetrada-umbral-k}
\end{equation}
This tetrad is obtained after constructing \(K\). In this unit it is not yet
called a Witt star: that object will be defined only after
constructing the ternary code, its hexads, and the Steiner system.

\subsection{Causal independence with respect to the physical output}

\begin{theorem}[Isolation of \(U\) and \(K\)]
\label{p3-83-c25-s4:thm:aislamiento-u-k-alpha-codata}
The composition
\begin{equation}
 \Omega_{\rm seed}\longrightarrow x_{\rm term}
 \xrightarrow{R_{\rm sgn}}U
 \xrightarrow{\mathcal H_{12}^{-1}}K
 \label{p3-83-c25-s4:eq:composicion-upstream-u-k}
\end{equation}
does not receive \(\alpha\), \(\alpha^{-1}\), CODATA or a target decimal
string as input. The seal \(K\) is determined before defining the
recurrence that will produce \(\alpha\).
\end{theorem}

\begin{proof}
The first part of the composition uses only seeds, emissions,
ternary readers, \(L_0,L_1\) operators, survival relations,
phase transport and the charts of the TPK state. The reading
\(R_{\rm sgn}=\Pi_H\circ E_{108}^{90,120}\circ R_{12}\) is the causal
composition of the terminal register, and the last arrow is the linear transformation
\(\mathcal H_{12}^{-1}=\mathcal H_{12}/4\)
on its integral image. None of the formulas contains a physical
variable or a metrological table. The base-one-thousand recurrence is introduced only
in the next unit, after \(K\) has been fixed.
\end{proof}

\subsection{Obstructions of the dodecaphasic closure: seal, carry, reversibility and non-circularity}

\begin{remark}[Certificate and validity control: Refutation of the record and the seal]
The construction is refuted if any of the following facts is verified:
\begin{enumerate}
 \item the windows \(E_m\) do not partition the one hundred eight steps;
 \item the chart \(\mathcal L_{12}\) fails its inverse or its congruences;
 \item the terminal census does not produce \(331\) triples, two matrices and a unique
       axial charge;
 \item an attempt is made to obtain \(U\) from \(B_{\rm term}\) after forgetting the
       signed coordinates;
 \item the naturality square
       \eqref{p3-83-c25-s4:eq:cuadrado-naturalidad-lector-firmado} fails;
 \item \(H_4^2\neq4I_4\), some preimage is not integral or the
       reconstructed vector differs from \eqref{p3-83-c25-s4:eq:sello-k-doce-autonomo};
 \item the Smith normal form differs from
       \eqref{p3-83-c25-s4:eq:snf-hadamard-autonoma};
 \item the certificate \((D_3,D_4,Q)\) loses rank, does not match the
       ledger’s prior emissions or does not reproduce \(U\);
 \item \(B_K\) is introduced before constructing \(K\), or is assigned a
       Witt structure before defining \(W_{12}\);
 \item a digit of \(\alpha\) or a metrological reference intervenes in
       any of the arrows of \eqref{p3-83-c25-s4:eq:composicion-upstream-u-k}.
\end{enumerate}
\end{remark}


% Unidad U017. Cierre afín, prolongación y semántica de frontera de alpha.
% Redacción autónoma desde la autoridad material 1063 y su sucesor V2.

\section{Coinductive prolongation of the dodecaphase cocycle and compatibility of the boundary \texorpdfstring{\(662/663\)}{662/663}}
\label{p3-83-c25-s5:chap:acarreo-alpha-frontera-662-663}

\subsection{Aritmetic closure domain}

The construction receives four dodecaphasic coordinates obtained in the preceding units:
\[
\begin{aligned}
 P={}&(141,592,653,589,793,238,462,643,383,279,502,884),\\
 E={}&(718,281,828,459,045,235,360,287,471,352,662,497),\\
 \Phi={}&(618,033,988,749,894,848,204,586,834,365,638,117),\\
 K={}&(234,543,140,729,659,824,621,058,914,794,146,601).
\end{aligned}
\]
The first three vectors are the first twelve triads of the fractional
Archimedean readings of \(\pi\), \(e\), and \(\varphi\). The fourth is
the dodecaphase coordinate of the enriched terminal state of the Topological
Phase Kernel, constructed before the carry operation is opened.

We fix the base.
\[
 B:=10^3=1000
\]
and the concatenation
\begin{equation}
 \iota(v_1,\ldots,v_{12})
 :=\sum_{m=1}^{12}v_mB^{12-m}.
 \label{p3-83-c25-s5:eq:concatenacion-base-mil-alpha}
\end{equation}
\Needspace{8\baselineskip}
In particular,
\[
\begin{aligned}
 \iota(P)&=141592653589793238462643383279502884,\\
 \iota(E)&=718281828459045235360287471352662497,\\
 \iota(\Phi)&=618033988749894848204586834365638117,\\
 \iota(K)&=234543140729659824621058914794146601.
\end{aligned}
\]

The operation studied in this unit is the map
\begin{equation}
 \mathfrak C_{12}:
 (P,E,\Phi,K,c_{13})
 \longmapsto(A,C),
 \label{p3-83-c25-s5:eq:mapa-cierre-dodecafasico-alpha}
\end{equation}
where \(A=(A_1,\ldots,A_{12})\) is the output word and
\[
 C=(c_1,\ldots,c_{13})
\]
is the complete system of thirteen Euclidean quotients: twelve internal
quotients and one right boundary condition.

\subsection{Left-to-right oriented Euclidean division}

For \(m=12,11,\ldots,1\), define
\begin{equation}
\begin{aligned}
 s_m&:=P_m+E_m-\Phi_m-K_m+c_{m+1},\\
 A_m&:=s_m\bmod B,
       \qquad 0\leq A_m<B,\\
 c_m&:=\frac{s_m-A_m}{B}
      =\left\lfloor\frac{s_m}{B}\right\rfloor .
\end{aligned}
\label{p3-83-c25-s5:eq:recurrencia-corta-alpha}
\end{equation}
\Needspace{9\baselineskip}
The second equality for \(c_m\) is also valid when \(s_m<0\),
because Euclidean division with remainder in
\(\{0,\ldots,B-1\}\) is used. Therefore,
\[
 -431=569-1000,
 \qquad
 -717=283-1000,
 \qquad
 -1200=800-2\cdot1000,
\]
and negative loans do not introduce ambiguity.

\begin{theorem}[Existence and Uniqueness of the Oriented Carry]
\label{p3-83-c25-s5:thm:unicidad-acarreo-orientado-alpha}
Given \(P,E,\Phi,K\) and the terminal carry \(c_{13}\), the recurrence
\eqref{p3-83-c25-s5:eq:recurrencia-corta-alpha} uniquely determines
\[
 (A_1,\ldots,A_{12};c_1,\ldots,c_{12}).
\]
\end{theorem}

\begin{proof}
Row \(m=12\) is determined by the Euclidean division of
\[
 P_{12}+E_{12}-\Phi_{12}-K_{12}+c_{13}
\]
by \(B\). Once \(c_{12}\) is known, row \(m=11\) is determined by
the same property. Descending induction reaches \(m=1\). At each
step, uniqueness of the Euclidean quotient and remainder prevents two distinct
outputs.
\end{proof}

\subsection{Finite dodecaphasic window}

The finite window is defined by the terminal condition
\[
 c_{13}^{\rm fin}=0.
\]
The complete evaluation is as follows. The table is printed in ascending
order, although it must be read causally from the last row to the
first.

\begingroup
\scriptsize
\setlength{\tabcolsep}{3.2pt}
\begin{longtable}{@{}r|rrrr|r|r|r|r@{}}
\caption{The twelve Euclidean divisions and thirteen carries of the finite
dodecaphasic window.}
\label{p3-83-c25-s5:tab:acarreo-alpha-finito-autonomo}\\
\toprule
\(m\)&\(P_m\)&\(E_m\)&\(\Phi_m\)&\(K_m\)&
\(c_{m+1}\)&\(s_m\)&\(A_m\)&\(c_m\)\\
\midrule
\endfirsthead
\toprule
\(m\)&\(P_m\)&\(E_m\)&\(\Phi_m\)&\(K_m\)&
\(c_{m+1}\)&\(s_m\)&\(A_m\)&\(c_m\)\\
\midrule
\endhead
1  &141&718&618&234& \(0\)& \(7\)&007& \(0\)\\
2  &592&281&033&543& \(0\)& \(297\)&297& \(0\)\\
3  &653&828&988&140& \(-1\)& \(352\)&352& \(0\)\\
4  &589&459&749&729& \(-1\)& \(-431\)&569& \(-1\)\\
5  &793&045&894&659& \(-2\)& \(-717\)&283& \(-1\)\\
6  &238&235&848&824& \(-1\)& \(-1200\)&800& \(-2\)\\
7  &462&360&204&621& \(0\)& \(-3\)&997& \(-1\)\\
8  &643&287&586&058& \(-1\)& \(285\)&285& \(0\)\\
9  &383&471&834&914& \(-1\)& \(-895\)&105& \(-1\)\\
10 &279&352&365&794& \(0\)& \(-528\)&472& \(-1\)\\
11 &502&662&638&146& \(0\)& \(380\)&380& \(0\)\\
12 &884&497&117&601& \(0\)& \(663\)&663& \(0\)\\
\bottomrule
\end{longtable}
\endgroup

The complete word of thirteen carries is
\[
 \boxed{
 C^{\rm fin}
 =(c_1,\ldots,c_{13})
 =(0,0,0,-1,-1,-2,-1,0,-1,-1,0,0,0).}
\]
The dodecaphasic output is
\[
 \boxed{
 A^{\rm fin}
 =(007,297,352,569,283,800,997,285,105,472,380,663).}
\]
We define the finite rational
\begin{equation}
 \boxed{
 \alpha_{12}^{\rm fin}
 :=10^{-36}\iota(A^{\rm fin})
 =\frac{7297352569283800997285105472380663}{10^{36}}.}
 \label{p3-83-c25-s5:eq:alpha-finita-autonoma}
\end{equation}
Therefore,
\[
 \alpha_{12}^{\rm fin}
 =0.007297352569283800997285105472380663.
\]

\subsection{Posterior metrological recognition: CODATA 2018}

The metrological comparison is performed only after closing the recurrence,
the right boundary, and the preceding rational. The integer identity gives
\begin{equation}
 \boxed{
 \alpha_{12}^{\rm fin}
 =10^{-36}
 \left\lfloor\frac{10^{36}}{137.035999084}\right\rfloor.}
 \label{p3-83-c25-s5:eq:alpha-finita-codata-2018}
\end{equation}
Consequently,
\begin{equation}
 (\alpha_{12}^{\rm fin})^{-1}
 =137.03599908400000000000000000000001066588\ldots,
 \label{p3-83-c25-s5:eq:alpha-inversa-codata-2018}
\end{equation}
and exactly reproduces the central chain published by CODATA 2018
\cite{Tiesinga2021CODATA}:
\[
 \boxed{137.035999084.}
\]
The word \emph{exactly} refers to the published central chain, not to a purported experimental measurement of thirty-six digits. The subsequent digits in \eqref{p3-83-c25-s5:eq:alpha-inversa-codata-2018} are consequences of the HMT rational already fixed. CODATA does not select the state, does not determine (K), does not decide the direction of the carry, and does not intervene in the dodecaphase table.

\subsection{Complete affine identity}

Each row of the table satisfies exactly
\begin{equation}
 P_m+E_m-\Phi_m-K_m+c_{m+1}
 =A_m+Bc_m.
 \label{p3-83-c25-s5:eq:identidad-local-acarreo-alpha}
\end{equation}
Introducing
\[
 C^-:=(c_1,\ldots,c_{12}),
 \qquad
 C^+:=(c_2,\ldots,c_{13}),
\]
the twelve local identities come together in
\begin{equation}
 \boxed{
 P+E-\Phi-K+C^+-BC^-=A.}
 \label{p3-83-c25-s5:eq:identidad-afin-vectorial-alpha}
\end{equation}
The term
\[
 \partial_BC:=C^+-BC^-
\]
is the carry cocycle of the oriented chain. Consequently,
\(P+E-\Phi-K=A\) is not a coordinatewise equality in
\(\mathbb Z^{12}\): it becomes correct only upon incorporating
\(\partial_BC\).

\begin{theorem}[Telescoping of the Thirteen Carries]
\label{p3-83-c25-s5:thm:telescopado-acarreos-alpha}
For any terminal condition \(c_{13}\), the recurrence satisfies
\begin{equation}
 \boxed{
 \iota(P)+\iota(E)-\iota(\Phi)-\iota(K)-\iota(A)
 =B^{12}c_1-c_{13}.}
 \label{p3-83-c25-s5:eq:telescopado-general-alpha}
\end{equation}
In the finite window, \(c_1=c_{13}=0\), and therefore
\begin{equation}
 \boxed{
 \iota(P)+\iota(E)-\iota(\Phi)-\iota(K)
 =\iota(A^{\rm fin}).}
 \label{p3-83-c25-s5:eq:identidad-entera-alpha-finita}
\end{equation}
\end{theorem}

\begin{proof}
We multiply \eqref{p3-83-c25-s5:eq:identidad-local-acarreo-alpha} by
\(B^{12-m}\) and sum over \(m=1,\ldots,12\). The contribution of the carries is
\[
 \sum_{m=1}^{12}
 (Bc_m-c_{m+1})B^{12-m}.
\]
The internal components cancel:
\begin{align*}
 \sum_{m=1}^{12}(Bc_m-c_{m+1})B^{12-m}
 &=\sum_{m=1}^{12}c_mB^{13-m}
   -\sum_{m=1}^{12}c_{m+1}B^{12-m}\\
 &=B^{12}c_1-c_{13}.
\end{align*}
This proves \eqref{p3-83-c25-s5:eq:telescopado-general-alpha}. The specialization
\(c_1=c_{13}=0\) gives \eqref{p3-83-c25-s5:eq:identidad-entera-alpha-finita}.
\end{proof}

\subsection{Exact inversion of the seal \texorpdfstring{\(K\)}{K}}

The local recurrence is reversible. Solving for \(K_m\) in
\eqref{p3-83-c25-s5:eq:identidad-local-acarreo-alpha},
\begin{equation}
 \boxed{
 K_m=P_m+E_m-\Phi_m+c_{m+1}-A_m-Bc_m.}
 \label{p3-83-c25-s5:eq:inversion-local-k-desde-alpha}
\end{equation}
By concatenation,
\begin{equation}
 \boxed{
 \iota(K)
 =\iota(P)+\iota(E)-\iota(\Phi)-\iota(A)
  -B^{12}c_1+c_{13}.}
 \label{p3-83-c25-s5:eq:inversion-concatenada-k-desde-alpha}
\end{equation}
On the finite boundary,
\[
 \iota(K)
 =\iota(P)+\iota(E)-\iota(\Phi)-\iota(A^{\rm fin}).
\]

Algebraic reversibility does not invert the causal order. The effective order
is
\[
 x_{\rm term}
 \longrightarrow K
 \longrightarrow(P,E,\Phi,K,c_{13})
 \longrightarrow(A,C).
\]
Formula \eqref{p3-83-c25-s5:eq:inversion-local-k-desde-alpha} is a commutativity
check: it proves that the output and the carries recover the seal
that was already fixed. It is not a rule for retrospectively selecting
\(K\) from a desired value of \(\alpha\).

\subsection{Coinductive prolongation of the dodecaphasic seal and inverse carry recurrence in \(\mathcal C=\{-2,-1,0,1,2\}\)}

Let
\[
 (P_k)_{k\geq1},
 \qquad
 (E_k)_{k\geq1},
 \qquad
 (\Phi_k)_{k\geq1}
\]
the infinite sequences of triads produced by the three coinductive branches, and let
\[
 \kappa(k):=1+((k-1)\bmod12).
\]
The prolongation of the stamp is periodic:
\[
 K_k^{\rm per}:=K_{\kappa(k)}.
\]
For a finite depth \(N\) and a terminal condition
\(t=c_{N+1}\), the recurrence is
\begin{equation}
\begin{aligned}
 s_k^{(N,t)}
 &=P_k+E_k-\Phi_k-K_{\kappa(k)}+c_{k+1}^{(N,t)},\\
 A_k^{(N,t)}
 &=s_k^{(N,t)}\bmod1000,\\
 c_k^{(N,t)}
 &=\left\lfloor\frac{s_k^{(N,t)}}{1000}\right\rfloor,
 \qquad k=N,\ldots,1.
\end{aligned}
\label{p3-83-c25-s5:eq:recurrencia-larga-alpha}
\end{equation}

The set of terminal carries
\[
 \mathcal C:=\{-2,-1,0,1,2\}
\]
is invariant for the certified inputs: for each triad and each
\(c_{k+1}\in\mathcal C\), the quotient \(c_k\) again belongs to
\(\mathcal C\). No favorable terminal condition is chosen. All
five conditions are propagated, and only the prefix common to all of them is published.

The execution of certified depth employs
\[
\begin{array}{c|r}
\text{quantity}&\text{value}\\ \hline
\text{ternary blocks per channel}&3501\\
\text{trits per channel}&3501\cdot6=21006\\
\text{guard decimal digits}&10020\\
\text{guard triads}&3340\\
\text{propagated terminal conditions}&5\\
\text{common triads}&3339\\
\text{common digits}&10017\\
\text{published digits}&10000
\end{array}
\]
The five terminal branches coalesce before the published prefix. In
particular, no terminal condition extracted from a target value of
\(\alpha\) intervenes in the output.

\begin{theorem}[Stabilized Output Compatibility]
\label{p3-83-c25-s5:thm:compatibilidad-prolongacion-alpha}
Let \(A^{[N]}\) and \(A^{[N']}\), with \(N'>N\), be two runs whose
five terminal conditions have coalesced before position \(N\).
Then the truncation of \(A^{[N']}\) to its first \(N\) blocks
coincides with \(A^{[N]}\). The stabilized prefixes define a unique
section in the inverse limit of finite words.
\end{theorem}

\begin{proof}
The recurrence is deterministic from right to left once the carry entering a
position is known. Coalescence asserts that all remote
conditions produce the same carry before the boundary of the
prefix. From that point, the same triads \(P_k,E_k,\Phi_k\) and the same
periodic seal produce, by Euclidean uniqueness, the same remainders and
quotients at all preceding positions. This is exactly
commutativity with truncation.
\end{proof}

\subsection{Boundary originating from the remote tail}

The prolongation determines on the visible boundary.
\[
 c_{13}^{\infty}=-1.
\]
The twelve internal carries remain the same as in the finite window:
\[
 \boxed{
 C_{\leq13}^{\infty}
 =(0,0,0,-1,-1,-2,-1,0,-1,-1,0,0,-1).}
\]
Only the information entering from the thirteenth triad changes. The
two relevant divisions are
\[
\begin{aligned}
 s_{12}^{\infty}
 &=884+497-117-601-1=662,\\
 A_{12}^{\infty}&=662,
 \qquad c_{12}^{\infty}=0,
\end{aligned}
\]
and
\[
\begin{aligned}
 s_{13}^{\infty}
 &=197+757-720-234-1=-1,\\
 A_{13}^{\infty}&=999,
 \qquad c_{13}^{\infty}=-1.
\end{aligned}
\]
Therefore, the beginning of the prolonged word is
\[
 \boxed{
 A^{\infty}
 =(007,297,352,569,283,800,997,285,105,472,380,662,999,\ldots).}
\]
The corresponding expansion begins
\begin{equation}
 \boxed{
 \alpha_{\infty}
 =0.007297352569283800997285105472380662999564172539642\ldots.}
 \label{p3-83-c25-s5:eq:alpha-infinita-autonoma}
\end{equation}

The difference between both boundaries is local and exact:
\[
 c_{13}^{\rm fin}-c_{13}^{\infty}=1,
 \qquad
 A_{12}^{\rm fin}-A_{12}^{\infty}=1,
\]
while
\[
 A_m^{\rm fin}=A_m^\infty,
 \qquad 1\leq m\leq11.
\]

Applying the telescoping theorem to the first twelve blocks of the
prolongation,
\[
 \iota(P)+\iota(E)-\iota(\Phi)-\iota(K)
 -\iota(A^\infty_{\leq12})
 =-c_{13}^{\infty}=1.
\]
Then
\begin{equation}
 \boxed{
 \iota(A^{\rm fin})
 =\iota(A^\infty_{\leq12})+1.}
 \label{p3-83-c25-s5:eq:relacion-unidad-fronteras-alpha}
\end{equation}
The ending \(663\) does not contradict the ending \(662\): the former
is the finite representative obtained by closing the window with
\(c_{13}=0\); the latter retains the carry \(-1\) from the remote
tail.

\begin{figure}[htbp]
\centering
\begin{tikzpicture}[
  >=Stealth,
  node distance=12mm and 16mm,
  box/.style={draw,rounded corners=1.5mm,align=center,
    inner xsep=3.2mm,inner ysep=2.5mm},
  arr/.style={->,thick}
]
\node[box] (shared) {shared data\\
 \(P,E,\Phi,K\)};
\node[box,below left=of shared] (finite)
 {finite boundary\\\(c_{13}=0\)};
\node[box,below right=of shared] (remote)
 {remote tail\\\(c_{13}=-1\)};
\node[box,below=of finite] (a663)
 {\(A_{12}=663\)\\clausura finita};
\node[box,below=of remote] (a662)
 {\(A_{12}=662\), \(A_{13}=999\)\\prolonged prefix};
\draw[arr] (shared)--(finite);
\draw[arr] (shared)--(remote);
\draw[arr] (finite)--(a663);
\draw[arr] (remote)--(a662);
\draw[<->,densely dashed] (a663)--node[below,align=center]
 {exact integer difference \(1\)} (a662);
\end{tikzpicture}
\caption{Boundary bifurcation of a single recurrence. Two distinct generators
are not being compared: only the datum entering from the
right changes.}
\label{p3-83-c25-s5:fig:bifurcacion-frontera-alpha}
\end{figure}

\subsection{Truncamiento compatible and rewhereo certified to \(36\) digits}

For \(x\geq0\), define
\[
 \operatorname{Tr}_{36}(x)
 :=10^{-36}\left\lfloor10^{36}x\right\rfloor
\]
and
\[
 \operatorname{Rd}_{36}(x)
 :=10^{-36}\left\lfloor10^{36}x+\frac12\right\rfloor.
\]
Let
\[
 a:=7297352569283800997285105472380662.
\]
The expansion \eqref{p3-83-c25-s5:eq:alpha-infinita-autonoma} can be written as
\[
 \alpha_\infty=\frac{a+\rho}{10^{36}},
 \qquad
 \rho=0.999564172539642709\ldots,
\]
so that \(0.999<\rho<1\).

\begin{theorem}[Exact resolution of the boundary \(662/663\)]
\label{p3-83-c25-s5:thm:frontera-alpha-662-663}
The coinductive prolongation satisfies
\[
 \boxed{
 \operatorname{Tr}_{36}(\alpha_\infty)
 =0.007297352569283800997285105472380662,}
\]
while its rounding to thirty-six digits satisfies
\[
 \boxed{
 \operatorname{Rd}_{36}(\alpha_\infty)
 =0.007297352569283800997285105472380663
 =\alpha_{12}^{\rm fin}.}
\]
\end{theorem}

\begin{proof}
Since \(0\leq\rho<1\),
\[
 \left\lfloor10^{36}\alpha_\infty\right\rfloor=a.
\]
This gives the truncation ending in \(662\). Since \(\rho>1/2\),
\[
 \left\lfloor10^{36}\alpha_\infty+\frac12\right\rfloor=a+1,
\]
which ends in \(663\). By \eqref{p3-83-c25-s5:eq:alpha-finita-autonoma},
\((a+1)10^{-36}\) is exactly \(\alpha_{12}^{\rm fin}\).
\end{proof}

Moreover,
\[
 0<\alpha_{12}^{\rm fin}-\alpha_\infty
 =\frac{1-\rho}{10^{36}}<10^{-39}.
\]
The difference is not an arithmetic error in the table: it quantifies the
difference between two distinct boundary operators.

\subsection{Causal isolation from \texorpdfstring{\(\alpha\)}{alpha} and metrological data}

At depth \(N\), the full domain of the operation is
\[
 \mathscr D_N
 =\bigl((P_k)_{k=1}^{N},(E_k)_{k=1}^{N},
        (\Phi_k)_{k=1}^{N},K,c_{N+1}\bigr).
\]
The carry map
\[
 \mathfrak C_N:\mathscr D_N
 \longrightarrow
 \{0,\ldots,999\}^{N}\times\mathcal C^{N}
\]
is defined exclusively by \eqref{p3-83-c25-s5:eq:recurrencia-larga-alpha}. Its
domain contains no target value of \(\alpha\), its inverse, a metrological
table, a target decimal string or a terminal carry chosen by
comparison.

\begin{theorem}[No interference from the metrological target]
\label{p3-83-c25-s5:thm:no-interferencia-metrologica-alpha}
Let \(\mathscr Y\) be any space of external data containing metrological
values or comparison tables. The formal extension
\[
 \widetilde{\mathfrak C}_N:
 \mathscr D_N\times\mathscr Y
 \longrightarrow
 \{0,\ldots,999\}^{N}\times\mathcal C^{N},
 \qquad
 \widetilde{\mathfrak C}_N(d,y):=\mathfrak C_N(d),
\]
factorizes by the projection
\[
 \operatorname{pr}_{\mathscr D_N}:
 \mathscr D_N\times\mathscr Y\longrightarrow\mathscr D_N.
\]
Therefore, for any \(y,y'\in\mathscr Y\),
\[
 \widetilde{\mathfrak C}_N(d,y)
 =\widetilde{\mathfrak C}_N(d,y').
\]
No intervention on a value of \(\alpha\) or a metrological
table alters \(K\), the quotients, the remainders or the prefix produced.
\end{theorem}

\begin{proof}
The statement results from
\[
 \widetilde{\mathfrak C}_N
 =\mathfrak C_N\circ\operatorname{pr}_{\mathscr D_N}.
\]
All operations of \(\mathfrak C_N\) are additions, subtractions and Euclidean
divisions of the elements of \(\mathscr D_N\). No coordinate of
\(\mathscr Y\) appears in them.
\end{proof}

The complete causal order is
\[
 \text{enriched branches of }\pi,e,\varphi
 \longrightarrow(P,E,\Phi),
\]
\[
 x_{\rm term}\longrightarrow K,
\]
\[
 (P,E,\Phi,K,\text{boundary})
 \longrightarrow(A,C)
 \longrightarrow\alpha_\infty,
\]
\[
 \alpha_\infty
 \longrightarrow\text{external comparison}.
\]
The last arrow cannot be reversed to select any of the previous ones.

\subsection{Exact obstructions of the prolongation: boundary, truncation, and causal isolation}

\begin{remark}[Certificate and validity control: Refutation of the affine closure and its prolongation]
The block is refuted in its domain if any of the following events occurs:
\begin{enumerate}
 \item some row of the \cref{p3-83-c25-s5:tab:acarreo-alpha-finito-autonomo}
       fails to satisfy \(s_m=A_m+1000c_m\);
 \item the published vector of thirteen carries does not reproduce the twelve rows;
 \item the telescoping identity fails
       \[
       \iota(P)+\iota(E)-\iota(\Phi)-\iota(K)-\iota(A)
       =1000^{12}c_1-c_{13};
       \]
 \item the inversion \eqref{p3-83-c25-s5:eq:inversion-local-k-desde-alpha} does not recover
       the twelve coordinates of \(K\);
 \item remote propagation produces \(c_{13}\neq-1\) for the certified
       prefix;
 \item the five terminal conditions of
       \(\mathcal C=\{-2,-1,0,1,2\}\) do not coalesce before the published
       digits;
 \item the triad following \(662\) is not \(999\), in which case the
       identification of the rounding with \(663\) is no longer proved;
 \item the generator opens a target expansion of \(\alpha\), its inverse,
       a metrological table, or a terminal condition chosen to force
       the coincidence.
\end{enumerate}
\end{remark}

Thus, two distinct boundary identities are thereby demonstrated:
\[
 \boxed{
 663=\text{finite closure with }c_{13}=0
     =\text{rounding to 36 digits of }\alpha_\infty,}
\]
\[
 \boxed{
 662=\text{truncation to 36 digits of }\alpha_\infty
 \quad\text{with}\quad c_{13}=-1,\ A_{13}=999.}
\]
The central metrological result is finally published in the same residence as the derivation:
\begin{equation}
 \boxed{
 \begin{gathered}
  (\alpha_{12}^{\rm fin})^{-1}
  =137.03599908400000000000000000000001066588\ldots,\\
  \text{central CODATA 2018 string:}\qquad137.035999084.
 \end{gathered}}
 \label{p3-83-c25-s5:eq:resultado-final-alpha-codata-2018}
\end{equation}
The implication expresses a subsequent comparison with the central string published; it does not introduce that string into the constructor.
